% Auto-filled by scripts/phase2_cqr.py.
\begin{table}[t]
  \centering
  \caption{Phase-2 empirical coverage and mean interval width on the
    in-distribution test split ($N_{\mathrm{test}} = 500$, pooled across
    $\sim 1.56 \times 10^{6}$ nodes), for the raw Gaussian ensemble
    interval $\bar{\sigma} \pm z_{1-\alpha/2}\,s$ and the CQR-calibrated
    interval. $\hat{Q}$ is the conformal quantile correction learned on
    the validation split ($N_{\mathrm{val}} = 500$). Negative $\hat{Q}$
    means the Gaussian base interval was slightly over-wide and CQR
    tightens it; positive $\hat{Q}$ means CQR widens it to restore
    coverage.}
  \label{tab:phase2-coverage}
  \begin{tabular}{cccccc}
    \toprule
    Nominal [\%] &
    \multicolumn{2}{c}{Uncalibrated (Gaussian)} &
    \multicolumn{2}{c}{CQR-calibrated} &
    $\hat{Q}$ [MPa] \\
    \cmidrule(lr){2-3}\cmidrule(lr){4-5}
    & Coverage [\%] & Width [MPa] & Coverage [\%] & Width [MPa] & \\
    \midrule
    80 & 86.41 & 0.25 & 82.72 & 0.24 & -0.005 \\
85 & 89.44 & 0.28 & 87.03 & 0.27 & -0.004 \\
90 & 92.32 & 0.32 & 91.28 & 0.31 & -0.003 \\
95 & 95.19 & 0.38 & 95.53 & 0.38 & 0.001 \\

    \bottomrule
  \end{tabular}
\end{table}

\begin{table}[t]
  \centering
  \caption{Phase-2 conditional coverage at nominal $90\%$, sliced by
    quartile of each design parameter. Coverage is marginal across nodes
    within the quartile; width is the mean interval width (MPa) within
    the quartile. Coverage stays within $\pm 3$~pp of nominal across all
    12 quartiles, and width grows with the known stress-concentration
    geometry (smaller fillet $R$, farther hole $p$, narrower flange $W$
    produce tighter intervals on the bulk of the graph but comparable
    widths at the concentrator).}
  \label{tab:phase2-conditional}
  \begin{tabular}{llccc}
    \toprule
    Slice & Range & $n$ samples & Coverage [\%] & Mean width [MPa] \\
    \midrule
    R Q1 & [6.51, 7.43] & 125 & 91.89 & 0.26 \\
R Q2 & [7.43, 8.42] & 125 & 92.16 & 0.33 \\
R Q3 & [8.42, 9.26] & 125 & 91.78 & 0.31 \\
R Q4 & [9.26, 10.00] & 125 & 89.33 & 0.34 \\
p Q1 & [57.02, 61.83] & 125 & 91.44 & 0.20 \\
p Q2 & [61.83, 66.19] & 125 & 91.91 & 0.29 \\
p Q3 & [66.19, 69.08] & 125 & 91.01 & 0.33 \\
p Q4 & [69.08, 72.00] & 125 & 90.76 & 0.43 \\
W Q1 & [19.01, 20.92] & 125 & 90.80 & 0.20 \\
W Q2 & [20.92, 22.03] & 125 & 92.67 & 0.29 \\
W Q3 & [22.03, 23.00] & 125 & 91.71 & 0.34 \\
W Q4 & [23.00, 24.00] & 125 & 89.95 & 0.41 \\

    \bottomrule
  \end{tabular}
\end{table}
